
6.3.
Pour toute \(K\)-algèbre graduée \(A\), on désigne par \(A_{\mathbb Q}\) la \(K\)-algèbre \(A\otimes_{\mathbb Z}\mathbb Q\), graduée de façon évidente. D'après 1.13, on définit un homomorphisme fonctoriel : \(1+\widehat A^+\xrightarrow{\ \ch\ }\widehat A_{\mathbb Q}^{+}\) en donnant sa valeur pour \(1+X\), soit
\begin{equation}\tag{6.3.1}
  \ch(1+X)=\exp(X)-1.
\end{equation}

Cet homomorphisme est compatible avec la multiplication définie en 6.1 sur \(1+\widehat A^+\)~: il suffit d'après 1.13 de la montrer pour les éléments \(1+X,1+Y\) de \(1+\widehat{K[X,Y]}{}^{+}\)~; or on a :
\[
\begin{aligned}
 \ch\bigl((1+X)*(1+Y)\bigr)
 &= \ch\!\left(\frac{1+X+Y}{(1+X)(1+Y)}\right)\\
 &= \ch(1+X+Y)-\ch(1+X)-\ch(1+Y)\\
 &= \exp(X+Y)-\exp(X)-\exp(Y)+1\\
 &= \ch(1+X).\ch(1+Y).
\end{aligned}
\]

Notons \(\eta\) l'endomorphisme (fonctoriel par rapport à \(A\)) de \(\widehat A_{\mathbb Q}^{+}\), défini par :
\[
 \eta\!\left(\sum_{k\ge1} a_k t^k\right)
 =
 \sum_{k\ge1}\frac{(-1)^{k-1}}{(k-1)!}.a_k t^k.
\]
Alors, \(\forall\varphi\in1+\widehat A^+\),
\begin{equation}\tag{6.3.2}
  \ch(\varphi)=\eta(\Log\varphi).
\end{equation}
En effet, \(\ch\) et \(\eta\circ\Log\) sont deux homomorphismes fonctoriels de groupes abéliens, et il suffit pour qu'ils soient égaux qu'ils coïncident sur \(1+X\), ce qui est évident.

On en déduit que \(\ch\) commute aux opérations de \(K\)~: en effet, il suffit là encore de le montrer pour \(1+X\), et on a :
\[
\begin{aligned}
 \forall a\in K,\quad\ch(1+X)^a)
 &=\ch(\exp(a\Log(1+X)))\\
 &=\eta(a\Log(1+X))\\
 &=a.\eta(\Log(1+X))\\
 &=a.\ch(1+X).
\end{aligned}
\]
